\documentclass{article}
\usepackage[UTF8]{ctex}
\usepackage{amsmath,amssymb,amsfonts}
\usepackage{enumitem}
\usepackage{geometry}
\geometry{a4paper, margin=1in}
% === 错题本排版优化（不参与题目标识） ===
\setlength{\parindent}{0pt}
\setlength{\parskip}{0.35em plus 0.10em minus 0.08em}
\setlength{\emergencystretch}{3em}
\setlist{itemsep=0.25em, topsep=0.40em, parsep=0pt}
\allowdisplaybreaks[2]
% === 错题本排版优化结束 ===
\begin{document}
\section*{南开大学2025年高等代数真题}
\section*{49．南开大学2025年高等代数试卷}

\subsection{题目1}
1．（30分）若线性方程组

\[
\left\{\begin{array}{l}
x_{1}+x_{2}+a x_{3}=1 \\
x_{1}+a x_{2}+x_{3}=1 \\
a x_{1}+x_{2}+x_{3}=-2
\end{array}\right.
\]

有无穷多组解，求 \(a\) 的值，并写出方程组的一般解。

\subsection{题目2}
2．（20 分）计算行列式

\[
\left|\begin{array}{lllll}
0 & 1 & 2 & 3 & 4 \\
1 & 0 & 1 & 2 & 3 \\
2 & 1 & 0 & 1 & 2 \\
3 & 2 & 1 & 0 & 1 \\
4 & 3 & 2 & 1 & 0
\end{array}\right|
\]

\subsection{题目3}
3．（20 分）求矩阵

\[
A=\left(\begin{array}{cccc}
0 & 1 & 0 & 0 \\
-1 & 3 & 1 & 0 \\
1 & -2 & 0 & 0 \\
0 & 1 & 1 & 1
\end{array}\right)
\]

的若尔当标准形。

\subsection{题目4}
4．（20 分）若正定矩阵 \(A\) 满足

\[
A^{2}=\left(\begin{array}{cccc}
10 & -6 & 0 & 0 \\
-6 & 10 & 0 & 0 \\
0 & 0 & 10 & -6 \\
0 & 0 & -6 & 10
\end{array}\right)
\]

求所有这样的 \(A\) 。

\subsection{题目5}
5．（20 分）设 \(\mathcal{L}\) 是线性空间 \(V\) 上的线性变换，且 \(\operatorname{Ker} \mathcal{L}^{k}=\operatorname{Ker} \mathcal{L}^{k+1}\) 对某个正整数 \(k\) 成立，证明：

\[
\operatorname{Ker} \mathcal{L}^{j}=\operatorname{Ker} \mathcal{L}^{j+1}
\]

对任意整数 \(j \geq k\) 成立．

\subsection{题目6}
6．（20 分）设 \(A, B \in \mathbb{R}^{\mathrm{n} \times \mathrm{n}}\) 满足 \(A^{2}=A, B^{2}=B\) ，求证：\(A\) 与 \(B\) 相似当且仅当它们的秩相等．

\subsection{题目7}
7．（10 分）设 \(V^{*}\) 是线性空间 \(V\) 的对偶空间，\(V_{1}^{*}, V_{2}^{*}\) 是 \(V^{*}\) 的子空间．记

\[
\begin{gathered}
W=\left\{v \in V \mid f(v)=0, \forall f \in V_{1}^{*} \cap V_{2}^{*}\right\} \\
W_{i}=\left\{v \in V \mid f(v)=0, \forall f \in V_{i}^{*}\right\}, i=1,2 .
\end{gathered}
\]

证明：\(W=W_{1}+W_{2}\) 。

\subsection{题目8}
8．（10 分）设 \(B \in \mathbb{R}^{2 \times 2}, C \in \mathbb{R}^{2 \times n}, D \in \mathbb{R}^{n \times n}\) ，且 \(B, D\) 均为对称矩阵。设 \(B\) 的两个特征值为 \(\mu_{1}, \mu_{2}\) ，矩阵 \(A=\left(\begin{array}{cc}B & C \\ C^{\prime} & D\end{array}\right)\) 的特征值为 \(\lambda_{1}, \lambda_{2}, \cdots, \lambda_{n+2}\) ，求证：

\[
\min \left\{\lambda_{1}, \lambda_{2}, \cdots, \lambda_{n+2}\right\} \leq \min \left\{\mu_{1}, \mu_{2}\right\} .
\]
\end{document}
